\honest{Class Project}{Eliezer Yudkowsky}{2008}

This is a story. Jeffreyssai tells his students not to try to do as well as Einstein. Einstein ``spoke nonsense of an impersonal God,'' was ``too caught up in the drama'' of rejecting quantum mechanics to fix it, and, though he ``reasoned cleanly'' on General Relativity, took ten years. ``Too slow!'' \nb{The charges repeat Yudkowsky's own in ``My Childhood Role Model'' and ``Einstein's Speed.'' Einstein's path to General Relativity was not clean: from 1913 to 1915 he held a theory that was wrong.} ``You will aspire to do BETTER than Einstein or you may as well not bother!'' \nb{This is the thesis of ``Einstein's Superpowers,'' posted the day before.}

His students, he says, achieve great things because he expects much of them. Their assignment: the correct theory of quantum gravity, working on their own, in one month. ``We will see if you have learned speed.'' He leaves. \nb{In ``The Failures of Eld Science,'' the month was for a case where ``you already have enough experimental evidence to determine an answer.'' Quantum gravity has, as yet, no experimental tests.}

Styrlyn takes charge. Each student writes down angles alone, and no one may criticize yet. \nb{Research on brainstorming favors this order: people working separately produce more ideas than groups brainstorming together.} Brennan can think only of ``the obvious.''

Taji says to assume Eld science's avenues were blind alleys, and to stop if anything starts to look like string theory. Hiriwa answers, ``The opposite of folly is folly.'' Hiriwa's own angle is to get rid of the infinities. Yin found a Lojban inscription in an abandoned city: ``Eureka! Eliminate t from the equations.'' Styrlyn wants quantum physics in a spatially local form. \nb{Infinities and the role of time are recognized problems in quantum gravity. The posts linked for Yin's and Styrlyn's angles are not in the book.}

Brennan asks why the energy of the quantum Hamiltonian is ``necessarily the same'' as the energy that curves spacetime, as Einstein made inertial and gravitational mass one. Yin thinks it too obvious for Eld science to have missed. Taji suggests proving that an exception would break a conservation law. Hiriwa wants instead a view in which the two ``are one and cannot be divided even conceptually.'' \nb{In 1964 Steven Weinberg derived the equality of gravitational and inertial mass from quantum theory, a proof of the kind Taji proposes.}

After thirty-seven seconds of silence someone says, ``I have an idea\ldots'' I stop there, before any result.
